%! Piecewise & Step Functions — Unit 1, Lesson 1.4 — Guided Notes & Practice (Answer Key)
% Mirrors notes/main.tex EXACTLY: every \blank{} becomes a short \ans{}, every \vterm{} becomes a
% \vtermans{}{} of the same fixed row height, and every work block is authored byte-identically,
% so the two files paginate alike (the \newpage plan makes the layout deterministic). No teacher
% notes here — they live in the lesson plan.
\documentclass[12pt]{article}
\usepackage{algebra2-article}
\usepackage{algebra2-key}   % pulls in -boxes; do NOT also load -boxes
% Vocabulary rows: the key's counterpart of the blank's \vterm — same fixed height, the
% definition filled beside the term (two lines at most).
\newcommand{\vrowheight}{2.0cm}
\newcommand{\vtermans}[2]{%
  \par\noindent\begin{minipage}[t][\vrowheight][t]{\linewidth}%
    \vspace*{6pt}\textbf{\textcolor{forest}{#1:}}\quad\ans{#2}%
  \end{minipage}\par}
% Coordinate grid: {xmin}{xmax}{ymin}{ymax}
\newcommand{\lgrid}[4]{%
  \draw[step=1, linegray!40, very thin] (#1,#3) grid (#2,#4);
  \draw[->, semithick, charcoal] (#1,0)--(#2,0) node[below right, font=\scriptsize]{$x$};
  \draw[->, semithick, charcoal] (0,#3)--(0,#4) node[left, font=\scriptsize]{$y$};}
% Months grid for the streaming service: x = months 0..7, one y-unit = $12.
\newcommand{\mgrid}{%
  \draw[step=1, linegray!40, very thin] (0,0) grid (7,4);
  \draw[->, semithick, charcoal] (0,0)--(7.6,0) node[below right, font=\scriptsize]{$m$};
  \draw[->, semithick, charcoal] (0,0)--(0,4.6) node[left, font=\scriptsize]{\$};
  \foreach \x in {1,...,7}{\node[below, font=\tiny, charcoal] at (\x,0){$\x$};}
  \foreach \y/\lab in {1/12,2/24,3/36,4/48}{\node[left, font=\tiny, charcoal] at (0,\y){$\lab$};}}
% Endpoint dots: {color}{x}{y}. Closed = filled; open = white-filled ring.
\newcommand{\fdot}[3]{\fill[#1] (#2,#3) circle (3.2pt);}
\newcommand{\hdot}[3]{\draw[very thick, #1, fill=white] (#2,#3) circle (3.2pt);}

\begin{document}

\pageheader{Unit 1, Lesson 1.4}{Guided Notes \& Practice --- Answer Key}
% NAMESTRIP: no \namedateperiod here — mirrors the blank.

\vspace{6pt}

% ── PAGE 1: vocabulary, then the hook ────────────────────────────────────────
\begin{vocabbox}
\small
Fill in each term \textbf{as we name it} in the notes below.
\vtermans{Piecewise-defined function}{one function with different rules on different parts of its domain, written with a brace}
\vtermans{Boundary (breakpoint)}{an $x$-value where the rule changes; the piece whose inequality includes it ($\le,\ge$) owns it}
\vtermans{Open vs.\ closed endpoint}{open $\circ$ excludes the point ($<,>$); closed $\bullet$ includes it ($\le,\ge$)}
\vtermans{Continuous at a boundary / discontinuity (jump)}{the pieces meet if both rules give the same value there; otherwise the graph jumps}
\vtermans{Greatest-integer (step) function $\lfloor x\rfloor$}{the greatest integer $\le x$ --- rounds \emph{down}; constant on each $[n,n+1)$, then jumps}
\end{vocabbox}

\vspace{8pt}

\boxguard[12]
\begin{hookbox}
\small
A streaming service is \emph{free for the first 3 months}, then costs \emph{\$12 a month} for
every month after that. Let $C(m)$ be the \emph{total} you have paid after $m$ months.
\par\vspace{4pt}
Total paid after $2$ months: \ans{$0$} \quad after $5$ months: \ans{$24$} \quad
Could \emph{one} straight line describe the total for every month? \ans{no}
\par\vspace{5pt}
\textbf{Month 3 --- the last free month, or the first paid one?} Circle one: \quad
\textbf{free} \qquad \textbf{\$12} \qquad Why? \ans{free --- $3$ satisfies $0\le m\le3$ (take the vote; section 1 settles it)}
\par\vspace{4pt}
One function, two rules. Every question today is this picture: \textbf{which rule does this
month use --- and what does the graph do where the rule changes?}
\end{hookbox}

% ── PAGE 2 — I DO: section 1 ─────────────────────────────────────────────────
\newpage
\begin{notesbox}{1. One Function, Several Rules --- Evaluate by Choosing the Piece}
\small
A \textbf{piecewise-defined function} uses different rules on different parts of its domain,
written with a brace. Each rule comes with the inputs it is \emph{for}:
\[ C(m)=\begin{cases} 0, & 0\le m\le 3 \\ 12(m-3), & m>3 \end{cases} \]
\textbf{To evaluate: first pick the piece, then apply its rule.}
\begin{itemize}[leftmargin=1.5em, itemsep=2pt]
  \item $C(2)$: is $2\le3$? \ans{yes} \; so use the \ans{first} piece: $C(2)=$ \ans{$0$}
  \item $C(5)$: is $5>3$? \ans{yes} \; so use the second piece:
        \begin{work}
          C(5) &= 12(5-3) \\
               &= 24
        \end{work}
  \item $C(3)$: the \textbf{boundary}, where the rule changes. $3$ satisfies $0\le m\le3$ --- the
        $\le$ \emph{includes} it --- so $C(3)=$ \ans{$0$}: month $3$ is still free.
\end{itemize}
\textbf{Boundary rule:} an input on the boundary belongs to the piece whose inequality
\emph{includes} it --- the one with $\le$ or $\ge$, never the strict one. The story does not
decide it; the \ans{inequality} does.

\vspace{8pt}
\textbf{\textcolor{forest}{The V Is Two Pieces.}}
Lesson 1.3's V is two lines pasted together at the vertex, and now we can \emph{write} it that
way. When $x\ge0$ the bars change nothing; when $x<0$ they flip the sign:
\[ |x|=\begin{cases} \text{\ans{$x$}}, & x\ge0 \\[2pt] \text{\ans{$-x$}}, & x<0 \end{cases}
   \qquad\text{Check: } |-3|=-(-3)= \text{\ans{$3$}} \]
\textbf{The trail from Lesson 1.3,} $H(x)=2|x-3|+1$. The vertex $x=3$ is the boundary.
\begin{center}
\begin{minipage}[t]{0.47\linewidth}
For $x\ge3$, $|x-3|=x-3$:
\begin{work}
  H(x) &= 2(x-3) + 1 \\
       &= 2x - 5
\end{work}
\end{minipage}\hfill
\begin{minipage}[t]{0.47\linewidth}
For $x<3$, $|x-3|=-(x-3)=3-x$:
\begin{work}
  H(x) &= 2(3-x) + 1 \\
       &= 7 - 2x
\end{work}
\end{minipage}
\end{center}
The right ray has slope \ans{$2$} and the left ray has slope \ans{$-2$}. At the boundary
\emph{both} rules give $H(3)=$ \ans{$1$}, so the two pieces \ans{meet} (meet / jump) at
the vertex --- the graph is one connected V.
\end{notesbox}

% ── PAGE 3 — I DO: section 2 — the graph, and the crux ───────────────────────
\newpage
\begin{notesbox}{2. Reading the Graph: Open or Closed, Meets or Jumps}
\small
A \textbf{closed} dot $\bullet$ \emph{includes} its endpoint ($\le,\ge$); an \textbf{open} dot
$\circ$ \emph{excludes} it ($<,>$). Same story, two different functions --- the \emph{total} paid
so far, and the \emph{price this month}:
\[ C(m)=\begin{cases} 0, & 0\le m\le3 \\ 12(m-3), & m>3 \end{cases} \qquad\qquad
   P(m)=\begin{cases} 0, & 0\le m\le3 \\ 12, & m>3 \end{cases} \]
\vspace{-6pt}
\begin{center}
\renewcommand{\arraystretch}{1.3}
\begin{tabularx}{\linewidth}{@{}X|X@{}}
\textbf{Total paid $C(m)$} & \textbf{Price this month $P(m)$} \\
\hline
\centering\arraybackslash
\begin{tikzpicture}[scale=0.30, baseline=(current bounding box.center)]
  \mgrid
  \draw[very thick, forest] (0,0)--(3,0);
  \draw[very thick, forest] (3,0)--(7,4);
  \fdot{forest}{3}{0}
\end{tikzpicture}
&
\centering\arraybackslash
\begin{tikzpicture}[scale=0.30, baseline=(current bounding box.center)]
  \mgrid
  \draw[very thick, redacc] (0,0)--(3,0);
  \draw[very thick, redacc] (3,1)--(7,1);
  \fdot{redacc}{3}{0}
  \hdot{redacc}{3}{1}
\end{tikzpicture}
\\
At $m=3$ the first rule gives \ans{$0$}; just past $3$ the second gives \ans{$0$} &
At $m=3$ the first rule gives \ans{$0$}; just past $3$ the second gives \ans{$12$} \\
Same value? \ans{yes} \; The pieces \ans{meet} (meet / jump) & Same value? \ans{no} \; The pieces \ans{jump} (meet / jump) \\
\emph{Continuous} at $3$: constant on \ans{$[0,3]$}, increasing on \ans{$[3,\infty)$} & A \emph{discontinuity} at $3$; the range is only \ans{$\{0,12\}$} \\
\end{tabularx}
\end{center}
\textbf{Meets or jumps --- one test:} both rules agree at the boundary, or they do not.

\vspace{6pt}
\textbf{\textcolor{forest}{The Crux: Exactly One Piece Owns the Boundary.}}
\begin{center}
\fcolorbox{goldacc}{goldbg}{\parbox{0.94\linewidth}{\small
A classmate reads $P(3)=12$ ``because the price is \$12 after the free months.'' But $3$
satisfies $0\le m\le3$, so $P(3)=$ \ans{$0$}, and the \ans{closed} (open / closed) dot at
$(3,0)$ says so. The open dot at $(3,12)$ is \emph{not} a point of the graph --- it marks where
the second piece starts \emph{without} its endpoint. Give a boundary to ``both'' pieces, or to
the strict one, and you have read the wrong dot.}}
\end{center}

\vspace{6pt}
\textbf{\textcolor{forest}{Step Functions: Rounding Down.}}
$P(m)$ is \emph{constant} on each piece, then jumps --- a \textbf{step function}. The most famous
one is the \textbf{greatest-integer function} $\lfloor x\rfloor$: the greatest integer
\emph{less than or equal to} $x$. It \textbf{rounds down}.
\begin{center}
\begin{minipage}[c]{0.58\linewidth}
$\lfloor 3.7\rfloor=$ \ans{$3$} \quad $\lfloor 5\rfloor=$ \ans{$5$} \quad
$\lfloor -0.4\rfloor=$ \ans{$-1$} \quad $\lfloor -2.5\rfloor=$ \ans{$-3$} \\[4pt]
\emph{Careful:} down means \emph{more negative}, not ``toward zero'' --- $-2.5$ goes to
\ans{$-3$}, not $-2$. \\[4pt]
The staircase: on $[n,n+1)$ the output is the constant $n$ --- closed on the left, open on the
right --- then it jumps by $1$. So $\lfloor x\rfloor=2$ for every $x$ in \ans{$[2,3)$}, and on
that interval the function is \ans{constant} (increasing / constant).
\end{minipage}\hfill
\begin{minipage}[c]{0.38\linewidth}\centering
\begin{tikzpicture}[scale=0.36]
  \lgrid{-1.5}{4.5}{-1.5}{4}
  \draw[very thick, forest] (-1,-1)--(0,-1); \fdot{forest}{-1}{-1} \hdot{forest}{0}{-1}
  \draw[very thick, forest] (0,0)--(1,0);    \fdot{forest}{0}{0}   \hdot{forest}{1}{0}
  \draw[very thick, forest] (1,1)--(2,1);    \fdot{forest}{1}{1}   \hdot{forest}{2}{1}
  \draw[very thick, forest] (2,2)--(3,2);    \fdot{forest}{2}{2}   \hdot{forest}{3}{2}
  \draw[very thick, forest] (3,3)--(4,3);    \fdot{forest}{3}{3}   \hdot{forest}{4}{3}
\end{tikzpicture}
\end{minipage}
\end{center}
\end{notesbox}

% ── PAGE 4 — WE DO: guided practice, alone on its page ───────────────────────
\newpage
\begin{practicebox}
\small
\textbf{We work this one together.} No story this time --- just the rule and its graph.
\[ p(x)=\begin{cases} -2x, & x\le -1 \\ x+3, & x>-1 \end{cases} \]
\vspace{-6pt}
\begin{center}
\begin{minipage}[c]{0.60\linewidth}
\begin{enumerate}[label=\textbf{\alph*.}, leftmargin=1.8em, itemsep=6pt]
  \item Pick the piece, then apply it. \\[3pt]
        $p(-3)=$ \ans{$6$} \quad $p(0)=$ \ans{$3$} \quad $p(2)=$ \ans{$5$} \\[3pt]
        $p(-1)=$ \ans{$2$} --- which piece owns $x=-1$? \ans{the first} \\[2pt]
        How do you know? \ans{$-1\le-1$: the $\le$ includes it; $-1>-1$ is false}
  \item \textbf{Meets or jumps at $x=-1$?} The first rule gives \ans{$2$}; just past
        $-1$ the second gives \ans{$2$}. So the graph \ans{meets} (meets / jumps):
        $p$ \emph{is} continuous there.
  \item From the graph: decreasing on \ans{$(-\infty,-1]$}, increasing on \ans{$[-1,\infty)$}; the minimum
        value is \ans{$2$}, at the boundary.
\end{enumerate}
\end{minipage}\hfill
\begin{minipage}[c]{0.36\linewidth}\centering
\begin{tikzpicture}[scale=0.34]
  \lgrid{-4.5}{4.5}{-1}{8}
  \draw[very thick, forest] (-4,8)--(-1,2);
  \draw[very thick, forest] (-1,2)--(4,7);
  \fdot{forest}{-1}{2}
\end{tikzpicture}
\end{minipage}
\end{center}
\begin{enumerate}[label=\textbf{\alph*.}, leftmargin=1.8em, itemsep=6pt, start=4]
  \item \textbf{The crux, on new ground.} A classmate says: ``Both rules give $2$ at $x=-1$, so it
        does not matter which piece owns it.'' Change the second piece to $x+4$ and test that.
        \\[3pt]
        Now which piece owns $x=-1$? \ans{still the first} \quad $p(-1)=$ \ans{$2$} \quad Just past
        $-1$ the second rule gives \ans{$3$} \\[3pt]
        So the graph now \ans{jumps} (meets / jumps) at $-1$ --- and the dot at $(-1,2)$ is
        \ans{closed} (open / closed) while the dot at $(-1,3)$ is \ans{open}. \\[3pt]
        What did the classmate miss? \ans{ownership is decided by the inequality, not by the rules happening to agree}
  \item \textbf{Round down.} $\lfloor 2.9\rfloor=$ \ans{$2$} \quad
        $\lfloor -1.2\rfloor=$ \ans{$-2$} \quad $\lfloor 4\rfloor=$ \ans{$4$} \quad
        $\lfloor x\rfloor=4$ for every $x$ in \ans{$[4,5)$} \\[3pt]
        Does the step $[4,5)$ own $x=5$? \ans{no} \; Which step does? \ans{$[5,6)$} \; So
        $\lfloor5\rfloor=$ \ans{$5$}
\end{enumerate}
\end{practicebox}

\end{document}
